\chapter{检索增强生成(RAG)}
\section{模型微调}
在前面的章节中，我们了解到大型语言模型（LLM）的强大能力。然而，LLM 也存在一些固有的局限性。首先是知识截止问题——模型的训练数据有截止日期，无法获取最新信息。
还有幻觉问题，即模型可能会生成看似合理但实际上错误的信息。此外，对于特定领域的私有数据，模型也无法直接访问，这限制了它在垂直领域的应用。

面对 LLM 的局限性，一种做法是对其进行微调（Fine-tuning）：在预训练模型的基础上，用领域特定的数据继续训练，让模型“学会”新的知识或任务。
技术上，微调就是在特定任务的数据集上，继续用梯度下降更新模型参数，最小化损失函数。
微调方法分为全量微调（Full Fine-Tuning）和参数高效微调（Parameter-Efficient Fine-Tuning, PEFT）。

\subsection{全量微调}
全量微调（Full Fine-Tuning）是最直接的微调方式：在预训练模型的基础上，用任务特定的数据集对所有模型参数进行更新。
全量微调的优点是效果好，模型可以充分学习任务特定的知识。但缺点也很明显：首先需要更新所有参数，对 GPU 显存和计算资源要求极高；
其次每个任务都需要保存一份完整的模型副本，存储开销大；此外还存在灾难性遗忘的风险：微调过程中可能丢失预训练阶段学到的通用知识。

\subsection{PEFT}
与全量微调不同，PEFT冻结预训练模型的大部分参数，只训练少量新增或特定的参数，从而大幅降低微调成本。
PEFT 方法可以大致分为三类：Adapter-based 方法、低秩方法和Prefix/Prompt 方法。

\subsubsection{Adapter-based 微调}
Adapter 方法由 Houlsby 等人于 2019 年提出，通过在预训练模型的 Transformer 层中插入小型的适配器模块（Adapter Module），
在微调过程中更新这个模块的参数来达到目的。

一个典型的 Adapter 模块包含：
\begin{enumerate}
	\item 一个降维投影层（Down-projection），将输入维度 $d$ 压缩到较小的维度 $r$
	\item 一个非线性激活函数（如 ReLU 或 GELU）
	\item 一个升维投影层（Up-projection），将维度从 $r$ 恢复回 $d$
\end{enumerate}

Adapter 的输出与原始层的输出通过残差连接相加：
\begin{equation*}
	h_{\text{out}} = h_{\text{in}} + \text{Adapter}(h_{\text{in}})
\end{equation*}

微调时，预训练模型的原始参数被冻结，只有 Adapter 模块的参数被更新。由于 $r \ll d$，Adapter 的参数数量远小于原始模型。
Adapter 方法具有模块化的特点：每个任务可以训练独立的 Adapter，推理时按需加载；多个任务的 Adapter 可以组合使用；不改变原始模型架构，易于部署。

\subsubsection{LoRA}
LoRA（Low-Rank Adaptation）由 Hu 等人于 2021 年提出，是一种基于低秩分解的高效微调方法。

LoRA 的核心假设是：模型在适应特定任务时，权重更新矩阵 $\Delta W$ 具有较低的“内在秩”。因此，可以用两个低秩矩阵 $A$ 和 $B$ 来近似 $\Delta W$：
\begin{equation*}
	W' = W + \Delta W = W + BA
\end{equation*}
其中 $W \in \mathbb{R}^{d \times k}$ 是预训练权重（冻结），$B \in \mathbb{R}^{d \times r}$，$A \in \mathbb{R}^{r \times k}$，$r \ll \min(d, k)$ 为秩超参数。

前向计算时：
\begin{equation*}
	h = Wx + BAx
\end{equation*}

LoRA 通常只应用于 Transformer 的注意力层（Query、Key、Value 投影矩阵），因为这些层被认为在任务适应中最为关键。

\subsubsection{Prefix 微调}
Prefix 微调（Prefix-Tuning）由 Li 和 Liang 于 2021 年提出，属于 Prompt-based 方法的一种。
与 Adapter 和 LoRA 在模型内部添加可训练参数不同，Prefix 微调在输入的 token 序列前面添加一组可学习的“前缀向量”。

具体来说，对于自回归语言模型，Prefix 微调在每一层 Transformer 的 Key 和 Value 缓存前面添加 $l$ 个可学习的虚拟 token：
\begin{equation*}
	\text{Attention}(Q, [P_K; K], [P_V; V])
\end{equation*}
其中 $P_K, P_V \in \mathbb{R}^{l \times d}$ 是可学习的前缀参数，$l$ 通常为 5-20。
微调时，预训练模型的所有参数保持冻结，只有前缀参数 $P_K$ 和 $P_V$ 被更新。

与 Prefix 微调相关的还有 Prompt Tuning（只学习输入 embedding 层的前缀）和 P-Tuning（在输入和中间层都添加前缀），
它们的核心思想相似，都是通过在输入端添加可学习的提示来引导模型行为。

\section{RAG}
微调的优点是模型将知识内化到参数中，推理时不需要额外的检索步骤，响应速度更快。但微调也有明显的缺点：需要高质量的训练数据，训练成本较高，而且一旦知识更新就需要重新训练。

与微调不同，检索增强生成（Retrieval-Augmented Generation，RAG）不需要修改模型参数，而是通过外挂知识库的方式来弥补 LLM 的知识缺陷。
它的做法是在生成回答之前，先从外部知识库中检索相关信息，然后将检索到的信息作为上下文提供给 LLM，让模型基于这些真实的信息生成回答。
这就像开卷考试，你不需要记住所有知识，只需要知道如何查找相关信息，然后基于这些信息回答问题。

\subsection{文本嵌入}
RAG 的第一步是将文本转换为向量表示，这个过程称为文本嵌入（Text Embedding）。
想象你在图书馆整理书籍。你会把相似主题的书放在同一个书架上。文本嵌入就是给每本书（文本）一个坐标（向量），使得内容相似的文本在向量空间中距离较近。

数学上，嵌入模型 $f$ 将文本 $t$ 映射到一个 $d$ 维向量：
\begin{equation*}
	f: t \rightarrow \mathbb{R}^d
\end{equation*}

例如：
\begin{verbatim}
文本1："机器学习是人工智能的一个分支"
文本2："深度学习属于机器学习领域"
文本3："今天天气很好"

嵌入后：
向量1：[0.2, 0.8, -0.3, ...]
向量2：[0.3, 0.7, -0.2, ...]  // 与向量1相似
向量3：[-0.5, 0.1, 0.9, ...]   // 与前两个不相似
\end{verbatim}

常用的相似度度量方法是余弦相似度，余弦值越接近 1，表示两个向量越相似。：
\begin{equation*}
	\text{similarity} = \cos(\theta) = \frac{\mathbf{A} \cdot \mathbf{B}}{\|\mathbf{A}\| \|\mathbf{B}\|}
\end{equation*}

有了文本与之对应的向量后，我们需要把他们存储下来，向量数据库就是这样一种专门用于存储和检索高维向量的数据库。
传统数据库就像电话簿，你通过名字（精确匹配）查找电话号码，
而向量数据库更像是一个智能地图，你“说我想找个附近好吃的中餐馆”，它会返回地理位置相近的选项。
常见的向量数据库有FAISS、Milvus和ChromaDB等。

\section{RAG 的优化}
之前介绍的是最基础的 RAG 实现方式，遵循"检索-阅读-生成"的简单流程：
\begin{enumerate}
	\item 将文档库切分为文本块并建立向量索引
	\item 将用户查询转换为向量，通过相似度搜索检索相关文本块
	\item 将检索结果与查询拼接为提示，输入 LLM 生成回答
\end{enumerate}

这种方式实现简单，但显然检索精度不高，很有可能检索到不相关的文档。
为此，研究者在基础 RAG 的基础上引入了多种优化策略，主要在检索前和检索后两个阶段进行改进。
检索前优化的目标是确保检索到更相关的文档，检索后优化则在检索完成后对结果进行精炼。

\subsection{检索前优化}
\subsubsection{查询改写}
查询改写（Query Rewriting）将用户的原始查询改写为更适合检索的形式。
例如，用户问“Python 怎么并行”，系统可以将其改写为“Python 并行计算 multiprocessing 多线程”，从而提高检索的准确性。

\subsubsection{查询路由}
查询路由（Query Routing）根据查询类型选择不同的检索策略或知识库。
例如，技术问题路由到技术文档库，而产品问题路由到产品手册库，避免在无关的知识库中浪费时间。

\subsubsection{查询分解}
查询分解（Query Decomposition）将复杂查询分解为多个子查询分别检索。
例如，用户问“对比 A 和 B 的优缺点”，系统可以分解为“列出 A 的优缺点”和“列出 B 的优缺点”两个子查询，分别检索后再合并结果。

\subsubsection{混合检索}
单一的向量检索或关键词检索各有优劣。向量检索擅长语义匹配，但可能忽略精确的关键词匹配；关键词检索（如 BM25）擅长精确匹配，但无法理解语义。

混合检索（Hybrid Search）将两种方法结合：
\begin{equation*}
	\text{score}(d, q) = \alpha \cdot \text{score}_{\text{dense}}(d, q) + (1 - \alpha) \cdot \text{score}_{\text{sparse}}(d, q)
\end{equation*}
其中 $\alpha$ 控制两种检索方式的权重，$\text{score}_{\text{dense}}$ 为向量检索的相似度分数，$\text{score}_{\text{sparse}}$ 为关键词检索的相关性分数。

\subsection{检索后优化}
\subsubsection{上下文压缩}
上下文压缩（Context Compression）从检索到的文档中提取与查询最相关的部分，减少噪声。
直接拼接整篇检索文档可能引入大量无关信息，压缩后只保留关键段落，既节省上下文窗口，又提高生成质量。

\subsubsection{上下文选择}
上下文选择（Context Selection）根据相关性阈值过滤低质量检索结果。
当检索返回的文档与查询相关性较低时，直接丢弃它们，避免不相关信息干扰 LLM 的生成。

\subsubsection{重排序（Reranking）}
初步检索返回的 top-$k$ 结果中，可能包含与查询不太相关的文档。重排序使用更精确的模型（通常是交叉编码器）对检索结果进行二次排序。

交叉编码器同时接收查询和文档作为输入，输出它们之间的相关性分数：
\begin{equation*}
	\text{score} = \text{CrossEncoder}(q, d)
\end{equation*}

与双编码器（分别编码查询和文档）相比，交叉编码器能捕捉查询和文档之间更细粒度的交互，但计算成本更高。因此，重排序通常只在初步检索的少量候选文档上进行。

\subsubsection{多跳检索}
对于需要多步推理的复杂问题，单次检索可能无法获取足够的信息。多跳检索（Multi-hop Retrieval）通过迭代检索来解决这个问题：
\begin{enumerate}
	\item 根据原始查询进行第一次检索
	\item 将检索结果与原始查询结合，生成新的查询
	\item 用新查询进行第二次检索
	\item 重复以上步骤，直到收集到足够的信息
\end{enumerate}

例如，对于问题“OpenAI 的 CEO 在哪个城市出生？”，系统可能需要先检索“OpenAI 的 CEO 是谁”，再检索“Sam Altman 在哪个城市出生”。

\subsection{高级 RAG 架构}
\subsubsection{Self-RAG}
Self-RAG 让 LLM 在生成过程中自主决定是否需要检索、何时检索，以及检索结果是否有用。
模型会生成特殊的“反思令牌”（reflection tokens）来指导检索行为，
包括判断当前是否需要检索外部信息的 \textbf{Retrieve}、评估检索到的文档是否与当前生成相关的 \textbf{IsRelevant}、
评估生成的内容是否有检索文档支持的 \textbf{IsSupported}，以及评估检索到的文档对回答是否有用的 \textbf{IsUseful}。
这种自适应的检索机制使得模型能够更灵活地利用外部知识，同时避免不必要的检索开销。

\subsubsection{Graph RAG}
传统的 RAG 基于文本块的向量检索，可能丢失文档之间的结构化关系。Graph RAG 将知识表示为知识图谱，
通过图结构来增强检索：从文档中提取实体和关系构建知识图谱，检索时不仅考虑文本相似度，还利用图结构进行多跳推理，通过图遍历发现实体之间的隐含关系。

Graph RAG 特别适合需要理解实体关系的复杂问答场景，如“某公司的 CEO 和 CTO 之间有什么共同的教育背景？”
